\subsection{The common abelian field at every dyadic exponent}

\begin{statement}{A.126}{Proved}
{Exact recursion and stabilization of the common dyadic fields}
Fix $D_1,\ldots,D_k\in\Q^\times$ and write
\[
 D_i=(-1)^{\delta_i}\prod_{P\in\mathcal P}P^{w_{iP}},
 \qquad \delta_i\in\{0,1\},\quad W=(w_{iP}).
\tag{A.126.1}
\]
For each $n=2^r$, $r\ge1$, let $B_n$ be the largest
abelian extension of $\Q$ contained in every field factor of
\[
 \mathcal E_n=\Q[Z_i^{\pm1}:1\le i\le k]/(Z_i^n-D_i)_i.
\tag{A.126.2}
\]
The square-coordinate maps give actual inclusions
$B_{n/2}\subset B_n$, and $B_2=\Q(\sqrt{D_i}:1\le i\le k)$.
For $n\ge4$ the following recursion computes the entire field.

When $\mu_n$ is not contained in $B_{n/2}$, one has
\[
             B_{2^s}=B_{n/2}\qquad(2^s\ge n).
\tag{A.126.3}
\]
When $\mu_n\subset B_{n/2}$, put $m=n/2$ and define
\[
 \begin{gathered}
 \mathcal R_m=\ker\bigl(W^{\mathsf t}:(\Z/m)^k
                                      \longrightarrow(\Z/m)^{\mathcal P}\bigr),\\
 0\le t_i<m,\qquad
 r_t=\prod_P P^{\frac1m\sum_i t_iw_{iP}}>0,
 \qquad \delta_t=\sum_i t_i\delta_i\pmod2,\\
 \prod_iD_i^{t_i}=(-1)^{\delta_t}r_t^m.
 \end{gathered}
\tag{A.126.4}
\]
Let $\chi_{2n}$ be the even character of
$(\Z/2n)^\times=\langle-1\rangle\times\langle5\rangle$
with $\chi_{2n}(5)=1/m$.
For each relation set
\[
 \lambda_t=\chi_{r_t}+\delta_t\chi_{2n},
 \qquad K_t=\overline\Q^{\ker\lambda_t}.
\tag{A.126.5}
\]
Here $\chi_{r_t}$ is the quadratic character of $r_t$;
a trivial character defines $K_t=\Q$.
Then
\[
                    B_n=B_{n/2}(K_t:t\in\mathcal R_m).
\tag{A.126.6}
\]
Every displayed field is included in every actual factor
by the root maps in the proof.  The relative extension
$B_n/B_{n/2}$ is elementary abelian of exponent two.

There is also an exact alternative for the infinite tower.
An integral vector $t$ with
\[
                 W^{\mathsf t}t=0,\qquad
                         \sum_i\delta_i t_i=1\pmod2
\tag{A.126.7}
\]
gives $\prod_iD_i^{t_i}=-1$ and forces
$\Q(\mu_{2n})\subset B_n$ for every dyadic $n$.
In the absence of such a vector, let $d_1,\ldots,d_\rho$
be the nonzero integral Smith invariants of $W^{\mathsf t}$,
and set
\[
       a=\max_j v_2(d_j),\qquad a=0\text{ when }\rho=0.
\tag{A.126.8}
\]
The tower is then constant for every $n\ge2^{a+2}$.
For $a=0$ it is already constantly $B_2$.
Thus (A.126.7) decides exactly whether this common-field
tower is unbounded.  For any odd positive integer $v$ and
any dyadic $n\ge4$, the same radicands also satisfy
\[
                            B_{nv}=B_n.
\tag{A.126.21}
\]
Here $B_{nv}$ is defined by (A.126.2) with exponent $nv$;
the equality uses the actual map $z_i\mapsto z_i^v$.
This statement concerns the actual
finite fibre fields and their inclusions; no lifting of
sections or compatibility of new rational coefficients is
asserted by it.
\end{statement}
\begin{proof}
\emph{One finite ambient field and all root stabilizers.}
For a dyadic integer $N\ge4$ put
\[
 F_N=\Q(\mu_{2N},\sqrt P:P\in\mathcal P),\qquad
 L_N=F_N(P^{1/N}:P\in\mathcal P),\qquad
 \Gamma_N=\operatorname{Gal}(F_N/\Q).
\tag{A.126.9}
\]
The positive-root degree and commutator calculations of
A.57 apply literally to this prime basis.  They give
$\operatorname{Gal}(L_N/F_N)=(\Z/(N/2))^{\mathcal P}$
and identify this group with
$[\operatorname{Gal}(L_N/\Q),\operatorname{Gal}(L_N/\Q)]$.
For $\sigma\in\Gamma_N$ write
\[
 \sigma(\zeta_{2N})=\zeta_{2N}^{u},\qquad
 \sigma(\sqrt P)=(-1)^{\epsilon_P}\sqrt P.
\tag{A.126.10}
\]
These are actual compatible signs; in particular the
sign on $\sqrt2$, when present, is determined by $u$.
All lifts to $L_N$ send
$P^{1/N}$ to $\zeta_N^{\epsilon_P+2h_P}P^{1/N}$,
with $h_P\in\Z/(N/2)$.
The root tuples have coordinates
$z_i=\zeta_{2N}^{\delta_i+2e_i}\prod_P P^{w_{iP}/N}$.
A lift fixes one such tuple precisely when
\[
 (u-1)(\delta_i+2e_i)+2\sum_Pw_{iP}\epsilon_P
              +4\sum_Pw_{iP}h_P=0\pmod{2N}
                                      \quad(1\le i\le k).
\tag{A.126.11}
\]
Let $\mathcal A_N$ be the set of images of all these
stabilizers and $H_N=\langle\mathcal A_N\rangle$.
The commutator identity and fixed-field correspondence give
\[
                              B_N=F_N^{H_N}.
\tag{A.126.12}
\]
Every abelian field inside all factors is inside $F_N$,
and its fixing subgroup must contain every stabilizer
image.  Conversely a character of $F_N^{H_N}$ vanishes
on every stabilizer and hence occurs in every factor.
These two implications prove (A.126.12).

The algebra map $\mathcal E_{N/2}\to\mathcal E_N$
sending each lower coordinate to the square of the
corresponding upper coordinate is injective: over
$\overline\Q$ its map of root tuples is onto, with
each fibre of size $2^k$.
Each upper field factor therefore contains the selected
lower field factor, and hence its common abelian field.
At exponent two every factor is the field generated
by the square roots, with independent choices of their
signs giving its embeddings.  This proves the asserted
base case and all field inclusions.

\emph{The formula at a step containing its cyclotomic field.}
Take $N=n$ and assume $\mu_n\subset B_{n/2}$.
Every element of
$\Gamma'=\operatorname{Gal}(F_n/B_{n/2})$
has $u\equiv1\pmod n$, and fixes the individual
square roots of the $D_i$.
In particular $\sum_Pw_{iP}\epsilon_P$ is even.
Set
\[
             A_i(\sigma)=\frac12\sum_Pw_{iP}\epsilon_P
                              +\frac{u-1}{4}\delta_i.
\tag{A.126.13}
\]
Dividing (A.126.11) by four gives
$Wh+(u-1)e/2=-A\pmod m$.
Here $(u-1)/2$ is zero modulo $m$, so the exact
condition for an actual stabilizer lift is
\[
                             Wh=-A\pmod m.
\tag{A.126.14}
\]
The perfect dot pairing on $(\Z/m)^k$ identifies
the annihilator of $\operatorname{im}W$ with
$\ker W^{\mathsf t}$.
Explicitly, integral unimodular row and column operations
put $W$ in Smith form.  A scalar equation $dh=b\pmod m$
is soluble exactly when $\gcd(d,m)$ divides $b$;
the vectors annihilating multiplication by $d$ impose
exactly this divisibility.  The zero rows require the
corresponding entries of $b$ to vanish.  Transporting
these pairings back gives the stated annihilator equality.
Consequently (A.126.14) is equivalent to
$t\cdot A=0\pmod m$ for every $t\in\mathcal R_m$.

For the representatives in (A.126.4), direct substitution
into (A.126.13) gives, in $\Q/\Z$,
\[
 \frac{t\cdot A}{m}
   =\frac12\sum_Pv_P(r_t)\epsilon_P
       +\frac{u-1}{2n}\sum_i t_i\delta_i
   =\chi_{r_t}(\sigma)+\delta_t\chi_{2n}(\sigma).
\tag{A.126.15}
\]
The final equality uses only $u=1$ or $1+n$ modulo
$2n$.  The identity
$5^{n/4}=1+n\pmod{2n}$ gives
$\chi_{2n}(1+n)=1/2$.
For $n=4$ it is $5=5\pmod8$; in higher degrees
successive squaring proves the displayed identity.
Thus the images of the actual stabilizers form the
subgroup
\[
              \Gamma'\cap\bigcap_{t\in\mathcal R_m}
                                      \ker\lambda_t.
\tag{A.126.16}
\]
Every stabilizer image is in $\Gamma'$ by the lower
coordinate map, and (A.126.14) constructs a lift and
tuple for every element of (A.126.16).
Taking fixed fields proves (A.126.6).
The characters in (A.126.15) restrict to values in
$\{0,1/2\}$ on $\Gamma'$, proving the relative
exponent assertion.

The field-factor embeddings are also explicit.
In any factor put $z_t=\prod_i z_i^{t_i}$.
For $\delta_t=0$ one has
$z_t=\zeta_n^j\sqrt{r_t}$.
Since this factor already contains $\mu_n$, it
contains $\sqrt{r_t}$.
For $\delta_t=1$ one has
$z_t=\zeta_{2n}^j\sqrt{r_t}$ with $j$ odd.
An automorphism fixing this root satisfies
$j(u-1)+n\epsilon=0\pmod{2n}$, where
$(-1)^\epsilon$ is its action on $\sqrt{r_t}$.
Hence $u=1+n\epsilon\pmod{2n}$, and
$\chi_{2n}+\chi_{r_t}$ vanishes on it.
The cyclic field $K_t$ is therefore in the actual
field $\Q(z_t)$ inside this factor.
This proves every asserted embedding without selecting
a preferred factor.

\emph{A missing cyclotomic layer stops all later steps.}
First suppose $i\notin B_2$.
For arbitrary $N\ge4$, let
$\Gamma''=\operatorname{Gal}(F_N/B_2)$.
For $\sigma\in\Gamma''$ with $u\equiv3\pmod4$,
the individual square-root condition is exactly the
divisibility by two required to solve (A.126.11)
for the $e_i$: $\gcd(u-1,N)=2$.
Thus every element in this cyclotomic coset belongs
to $\mathcal A_N$.
There is such an element $\tau$ because $i\notin B_2$.
Every element of $\Gamma''$ either lies in this
coset or becomes one after multiplication by $\tau$.
It follows that $H_N=\Gamma''$ and $B_N=B_2$.
This proves (A.126.3) when its first failure is $n=4$.

Now let $n\ge8$ be the first failed step.
All earlier steps used (A.126.6), and in particular
$\mu_{n/2}\subset B_{n/4}$.
Fix any dyadic $N\ge n$ and work in the single field
$F_N$.  Put
$\Gamma''=\operatorname{Gal}(F_N/B_{n/2})$.
Its cyclotomic values are $1$ modulo $n/2$.
Since $\mu_n\not\subset B_{n/2}$, its elements
with
\[
                         u=1+n/2\pmod n
\tag{A.126.17}
\]
form a nonempty coset of an index-two subgroup.
For each such element, $v_2(u-1)=\log_2(n)-1$.
In the divided equation (A.126.11), the coefficient
$(u-1)/2$ has image the multiples of $n/4$ modulo
$N/2$.  Therefore a lift fixing an $N$th-root tuple
exists exactly when
\[
                              Wh=-A\pmod{n/4}.
\tag{A.126.18}
\]
The formula proved at the preceding level $n/2$
identifies this condition exactly with fixing
$B_{n/2}$, for these elements fixing $\mu_{n/2}$.
This assertion follows from (A.126.14)--(A.126.16)
with $n$ replaced by $n/2$; reducing the present
$u,\epsilon,A$ gives those same equations.
Thus every element of the coset (A.126.17) belongs
to $\mathcal A_N$.
Choose one such $\tau$.  The same two-coset product
argument makes $\Gamma''\subset H_N$.
The lower coordinate maps give the reverse inclusion.
Hence $B_N=B_{n/2}$ for every $N\ge n$.
Together with the first case this proves permanent
stabilization at the first failed step, and therefore
the entire stated recursion.

\emph{The integral test and its finite stabilization bound.}
For a vector in (A.126.7), the invertible root product
$z_t=\prod_i z_i^{t_i}$ in any factor satisfies
$z_t^n=-1$.
For dyadic $n$ every such root is a primitive
$2n$th root of unity.
Thus $\Q(z_t)=\Q(\mu_{2n})$ is literally in every
factor.  These degrees are unbounded as $n$ increases.

In the absence of (A.126.7), the functional
$\delta:\Z^k\to\Z/2$ annihilates the integral
kernel of $W^{\mathsf t}$.
Smith form then proves the following precise claim:
\[
 W^{\mathsf t}t=0\pmod{2^{a+1}}
                         \quad\Longrightarrow\quad\delta(t)=0.
\tag{A.126.19}
\]
Indeed the nonzero Smith coordinates $d_jt'_j$
are divisible by $2^{a+1}$, so every corresponding
$t'_j$ is even.  The remaining coordinates lie
in the integral kernel.  Transporting back expresses
$t=2s+t_0$ with $W^{\mathsf t}t_0=0$; both summands
are killed by $\delta$.
For $a=0$, this says that no product of the $D_i$
can be a negative rational square.  The elementary
multiquadratic character pairing identifies
$i\in B_2$ exactly with such a product: its
squareclass is $-1$ in the span of the individual
squareclasses.  Therefore $i\notin B_2$, and the
first stabilization case gives $B_N=B_2$ for all $N$.
The rank-zero case has every $D_i=1$ under the
same assumption and also satisfies this conclusion.

Finally let $a\ge1$ and put $n_0=2^{a+2}$.
A failed step at or before $n_0$ has already proved stabilization.
Otherwise apply (A.126.6) at $n_0$.
Equation (A.126.19) makes every $\delta_t$ there zero,
so every added $K_t$ is quadratic.
The previous field $B_{n_0/2}$ is contained in
\[
                F_{n_0/2}=\Q(\mu_{n_0},\sqrt P:P\in\mathcal P).
\tag{A.126.20}
\]
Every character of this abelian field has order
dividing $2^a$; this includes $a=1$, when the
cyclotomic and quadratic factors all have exponent two.
The composite with the added quadratic fields has
the same exponent bound.
In contrast $\operatorname{Gal}(\Q(\mu_{2n_0})/\Q)$
has exponent $2^{a+1}$.
Containment would give a restriction surjection
from a group of exponent dividing $2^a$ onto this
larger-exponent group, which is impossible because
the image of every element is killed by $2^a$.
Thus $\mu_{2n_0}\not\subset B_{n_0}$.
The proved stopping rule at the next step shows
$B_N=B_{n_0}$ for every $N\ge n_0$.
This proves the bound and both directions of the
unboundedness test.

\emph{The exact odd-factor comparison.}
Fix dyadic $n\ge4$ and odd $v$ and put $N=nv$.
The same A.57 calculation for this exponent uses
$F_N=\Q(\mu_{2N},\sqrt P:P\in\mathcal P)$ and
$L_N=F_N(P^{1/N}:P\in\mathcal P)$.
Its abelianization is again $F_N$, and restriction
$\Gamma_N\to\Gamma_n$ is onto.
In (A.126.11), now modulo $2N$, write
$k_P=\epsilon_P+2h_P$.
The power map $z_i\mapsto z_i^v$ sends any fixed
$N$th-root tuple to a fixed $n$th-root tuple,
proving that its stabilizer image restricts to
$\mathcal A_n$.
Conversely take any lift $\sigma\in\Gamma_N$ of an
element of $\mathcal A_n$, and retain one actual
solution $(k_P,e_i)$ of its $n$th-root equations.
The Chinese remainder theorem gives new parameters
\[
 \begin{array}{lll}
 k'_P=k_P\pmod n,& k'_P=0\pmod v,\\
 e'_i=e_i\pmod n,& 2e'_i=-\delta_i\pmod v.
 \end{array}
\tag{A.126.22}
\]
Their parities are still the required $\epsilon_P$,
so they define actual lifts of this particular
$\sigma$ to $L_N$.
The root equations hold modulo $2n$ by the chosen
old solution.  Modulo $v$ both terms
$(u-1)(\delta_i+2e'_i)$ and $2\sum_Pw_{iP}k'_P$
vanish separately.  Since $\gcd(2n,v)=1$, the
equations hold modulo $2nv$.
This constructs an actual fixed upper tuple for
every such lift, and proves
$\mathcal A_N=\operatorname{res}^{-1}(\mathcal A_n)$.
The identity belongs to $\mathcal A_n$, so this
preimage contains the full restriction kernel.
Generating subgroups consequently gives
$H_N=\operatorname{res}^{-1}(H_n)$.
Their exact fixed fields prove (A.126.21).
\end{proof}

\noindent The exact finite certificate
\path{code/dyadic_common_field_tower_certificate.py}
compares the recursive fields with the groups generated
by the literal root-stabilizer equations at each tested
exponent.  It retains the compatible cyclotomic signs,
all relation vectors and the subgroup-generation step.
Its executed output is retained beside the source.
The all-exponent recursion and the stabilization bound
follow from the preceding proof.
